% ===========================================================================
% v0.23 — ΑΥΤΟΣ Ο ΚΩΔΙΚΑΣ ΕΙΝΑΙ ΚΟΙΝΟΣ ΣΤΙΣ ΔΥΟ ΕΚΔΟΣΕΙΣ
% SCREEN και PRINT διαφέρουν μόνο ως προς την παλέτα του βιβλίου.
% ===========================================================================
\thispagestyle{plain}
\begin{center}
  {\BookHeadingFont\Large\bfseries\color{BookPrimary}Δύο τύποι πινάκων}\\[3pt]
  {\small Ίδια τυπογραφία, διαστάσεις και κεφαλίδα — διαφορετικό φόντο γραμμών}
\end{center}
\vspace{6mm}
\BookMinorHeading{Α — Πίνακας με εναλλάξ χρωματισμό}
% Μόνο η λέξη BookStripedTable αλλάζει σε σχέση με τον τύπο Β.
\begin{BookStripedTable}
\begin{tabularx}{\linewidth}{@{}>{\raggedright\arraybackslash}X>{\centering\arraybackslash}p{25mm}>{\centering\arraybackslash}p{33mm}@{}}
\BookTableHeaderRow
\BookTableHead{Μέγεθος} & \BookTableHead{Σύμβολο} & \BookTableHead{Μονάδα S.I.}\\
Μήκος & $\ell$ & $\mathrm{m}$\\
Μάζα & $m$ & $\mathrm{kg}$\\
Χρόνος & $t$ & $\mathrm{s}$\\
Ταχύτητα & $\upsilon$ & $\mathrm{m/s}$\\
Επιτάχυνση & $a$ & $\mathrm{m/s^2}$\\
Δύναμη & $F$ & $\mathrm{N}$\\
\end{tabularx}
\end{BookStripedTable}
\vspace{9mm}
\BookMinorHeading{Β — Πίνακας με λευκές γραμμές}
% Τα ίδια δεδομένα, το ίδιο tabularx, η ίδια κεφαλίδα.
% Μόνο το BookPlainTable αντί BookStripedTable.
\begin{BookPlainTable}
\begin{tabularx}{\linewidth}{@{}>{\raggedright\arraybackslash}X>{\centering\arraybackslash}p{25mm}>{\centering\arraybackslash}p{33mm}@{}}
\BookTableHeaderRow
\BookTableHead{Μέγεθος} & \BookTableHead{Σύμβολο} & \BookTableHead{Μονάδα S.I.}\\
Μήκος & $\ell$ & $\mathrm{m}$\\
Μάζα & $m$ & $\mathrm{kg}$\\
Χρόνος & $t$ & $\mathrm{s}$\\
Ταχύτητα & $\upsilon$ & $\mathrm{m/s}$\\
Επιτάχυνση & $a$ & $\mathrm{m/s^2}$\\
Δύναμη & $F$ & $\mathrm{N}$\\
\end{tabularx}
\end{BookPlainTable}
\par\medskip\noindent
{\small\itshape\color{BookInk}Οι δύο πίνακες χρησιμοποιούν την ίδια κεφαλίδα και το ίδιο μέγεθος κειμένου. Δεν έχουν κάθετες γραμμές. Στο δίχρωμο PDF το αχνό φόντο του Α αποτελεί τονικότητα του δεύτερου μελανιού.}
